\section{Prime Factorization Changes}
\label{sec:factorization}

\subsection{Factorization Table}

\begin{table}[h]
\centering
\caption{Prime factorisation of 2020 and 2030 seat counts for states
         with projected changes. ``Depth'' is the depth of the \ar{}
         bisection tree.}
\label{tab:factorization}
\begin{tabular}{lccccc}
\toprule
State & 2020 $k$ & Factorisation & Depth & 2030 $k$ & Factorisation \\
\midrule
Texas & 38 & $2 \times 19$ & 2 & 41 & prime \\
California & 52 & $2^2 \times 13$ & 3 & 50 & $2 \times 5^2$ \\
Florida & 28 & $2^2 \times 7$ & 3 & 29 & prime \\
New York & 26 & $2 \times 13$ & 2 & 25 & $5^2$ \\
Arizona & 9 & $3^2$ & 2 & 10 & $2 \times 5$ \\
Georgia & 14 & $2 \times 7$ & 2 & 15 & $3 \times 5$ \\
N. Carolina & 14 & $2 \times 7$ & 2 & 15 & $3 \times 5$ \\
Illinois & 17 & prime & 1 & 16 & $2^4$ \\
Michigan & 13 & prime & 1 & 12 & $2^2 \times 3$ \\
Pennsylvania & 17 & prime & 1 & 16 & $2^4$ \\
\bottomrule
\end{tabular}
\end{table}

\subsection{Texas: The Most Dramatic Change}

Texas goes from $38 = 2 \times 19$ to $41$ (prime).
Under the current tree, Texas is bisected into two halves of 19 districts each;
each half is then recursively bisected 19 ways (19 is prime, so this is a
flat 19-way split).
The tree has depth 2 with two major subtrees.

Under the 2030 tree, Texas is directly split 41 ways at the root.
This is a flat tree of depth 1.
Every tract's district assignment is determined by a single 41-way
METIS partition rather than a hierarchical 2-level bisection.
The change from a 2-level to a 1-level tree means that the geographic
organization of the state's districts is fundamentally different.

\paragraph{Expected boundary movement.}
With $38 \to 41$ districts, the target district population decreases
from $P_{\rm state}/38$ to $P_{\rm state}/41$ (approximately 7.5\%
smaller per district).
Additionally, the tree structure changes.
We expect boundary movement of approximately 15--25\% of tracts
(compared to 3--5\% for seed changes within the current tree structure).

\subsection{California: Structural But Not Dramatic}

California goes from $52 = 2^2 \times 13$ to $50 = 2 \times 5^2$.
Both are composite, but the factorisation paths differ:
\begin{itemize}
\item $52$: bisect to 2 halves of 26 each; each 26 bisects to 2 halves of
      13 each; each 13 is a flat prime split.
      Tree depth: 3.
\item $50$: bisect to 2 halves of 25 each; each 25 splits as $5 \times 5$
      (bisect to 5 halves of 5 each; each 5 is a flat prime split).
      Tree depth: 3.
\end{itemize}

Both have depth 3, but the intermediate splitting differs.
At the second level, $52$-tree makes halves of 13 while $50$-tree makes
fifths of 5.
This is a structural change but less dramatic than the Texas prime transition.

\subsection{Illinois, Michigan, Pennsylvania: Prime to Composite}

These states currently have prime seat counts ($17, 13, 17$) and will
move to composite ($16, 12, 16$).
A prime $k$ has a flat single-level tree (all tracts assigned in one
$k$-way METIS partition).
A composite $k$ has a hierarchical tree.

This is the reverse of the Texas transition: going from flat to hierarchical.
Hierarchical trees are generally more compact (each subproblem is a cleaner
bisection) but the boundary movement from restructuring the tree is substantial.

\begin{table}[h]
\centering
\caption{Qualitative characterisation of the 2030 tree-structure changes.}
\label{tab:tree-changes}
\begin{tabular}{lcc}
\toprule
State & 2020 tree type & 2030 tree type \\
\midrule
Texas & Two-level (2-19) & Flat prime \\
California & Three-level ($4 \times 13$) & Three-level ($2 \times 25$) \\
Florida & Three-level ($4 \times 7$) & Flat prime \\
New York & Two-level (2-13) & Two-level (5-5) \\
Arizona & Two-level (3-3) & Two-level (2-5) \\
Georgia/NC & Two-level (2-7) & Two-level (3-5) \\
Illinois/PA & Flat prime (17) & Four-level ($2^4$) \\
Michigan & Flat prime (13) & Three-level ($4 \times 3$) \\
\bottomrule
\end{tabular}
\end{table}
